% =====================================================================
%  Preâmbulo base: pacotes, tipografia e convenções de todo documento
%  LaTeX do repositório
%
%  Origem: `config/preambulo.tex` do documento de projeto do rhae-sensor.
%  As seções abaixo seguem a ordem de lá, e o que difere está marcado
%  com DIFERENÇA.
%
%  Convenções tipográficas e de apresentação adotadas
%    ISO 216        formato A4
%    ISO 80000-1    números e unidades (vírgula decimal, espaço fino entre
%                   grupos de três algarismos, unidade em tipo romano)
%    Tabelas sem filetes verticais (booktabs); legenda de tabela acima,
%    de figura abaixo.
%
%  Uso, antes do `\begin{document}`:
%    report/*.tex          via `\usepackage{estilo-reporte}`, que carrega este
%    docs/product/*.tex    % =====================================================================
%  Preâmbulo base: pacotes, tipografia e convenções de todo documento
%  LaTeX do repositório
%
%  Origem: `config/preambulo.tex` do documento de projeto do rhae-sensor.
%  As seções abaixo seguem a ordem de lá, e o que difere está marcado
%  com DIFERENÇA.
%
%  Convenções tipográficas e de apresentação adotadas
%    ISO 216        formato A4
%    ISO 80000-1    números e unidades (vírgula decimal, espaço fino entre
%                   grupos de três algarismos, unidade em tipo romano)
%    Tabelas sem filetes verticais (booktabs); legenda de tabela acima,
%    de figura abaixo.
%
%  Uso, antes do `\begin{document}`:
%    report/*.tex          via `\usepackage{estilo-reporte}`, que carrega este
%    docs/product/*.tex    % =====================================================================
%  Preâmbulo base: pacotes, tipografia e convenções de todo documento
%  LaTeX do repositório
%
%  Origem: `config/preambulo.tex` do documento de projeto do rhae-sensor.
%  As seções abaixo seguem a ordem de lá, e o que difere está marcado
%  com DIFERENÇA.
%
%  Convenções tipográficas e de apresentação adotadas
%    ISO 216        formato A4
%    ISO 80000-1    números e unidades (vírgula decimal, espaço fino entre
%                   grupos de três algarismos, unidade em tipo romano)
%    Tabelas sem filetes verticais (booktabs); legenda de tabela acima,
%    de figura abaixo.
%
%  Uso, antes do `\begin{document}`:
%    report/*.tex          via `\usepackage{estilo-reporte}`, que carrega este
%    docs/product/*.tex    % =====================================================================
%  Preâmbulo base: pacotes, tipografia e convenções de todo documento
%  LaTeX do repositório
%
%  Origem: `config/preambulo.tex` do documento de projeto do rhae-sensor.
%  As seções abaixo seguem a ordem de lá, e o que difere está marcado
%  com DIFERENÇA.
%
%  Convenções tipográficas e de apresentação adotadas
%    ISO 216        formato A4
%    ISO 80000-1    números e unidades (vírgula decimal, espaço fino entre
%                   grupos de três algarismos, unidade em tipo romano)
%    Tabelas sem filetes verticais (booktabs); legenda de tabela acima,
%    de figura abaixo.
%
%  Uso, antes do `\begin{document}`:
%    report/*.tex          via `\usepackage{estilo-reporte}`, que carrega este
%    docs/product/*.tex    \input{../../report/config/preambulo}
%
%  FORMA DE MANUAL TÉCNICO. Fonte, números, legendas e pacotes vêm da
%  origem. A forma da página não: é de manual de referência, e não de
%  artigo. Coluna única, títulos recuados para a esquerda do texto,
%  rodapé com identificação e "página X de Y" em toda página, avisos com
%  palavra de sinal (ANSI Z535.6), passos numerados e histórico de
%  revisões (IEC/IEEE 82079-1). O desenho de página segue a classe
%  `refman` do CTAN.
%
%  O documento carrega `babel` ANTES deste arquivo. Documento deitado, que
%  existe para caber tabela larga, carrega também `geometry` antes; aí a
%  margem de manual não entra e os títulos não avançam.
%
%  Motor: pdflatex. Charis SIL, newtxmath e Bera Mono entram como Type 1
%  em T1, e a matemática do newtxmath não passa pelo `unicode-math`.
% =====================================================================
\RequirePackage{iftex}
\RequirePDFTeX
% Este arquivo entra tanto de um `.tex` quanto de dentro do
% `estilo-reporte.sty`, onde o `@` ja e letra. O codigo de categoria e
% devolvido como estava, e nao forcado a "outro" no fim.
\edef\pbrestauraarroba{\catcode`\noexpand\@=\the\catcode`\@\relax}
\makeatletter

% ---------------------------------------------------------------------
%  Fontes, idioma e microtipografia
% ---------------------------------------------------------------------
\usepackage[T1]{fontenc}
\usepackage[utf8]{inputenc}
% DIFERENÇA: o idioma é do documento, porque há documento em inglês.
\@ifpackageloaded{babel}{}{\usepackage[brazilian]{babel}}
\usepackage{CharisSIL}
\usepackage[charter]{newtxmath}
\usepackage[scaled=0.83]{beramono}
\usepackage[protrusion=true,expansion=true,tracking=true]{microtype}
\SetTracking{encoding=*,shape=sc}{40}
% DIFERENÇA: a origem não usa sem serifa. Os documentos daqui pedem
% `\sffamily` em título, tabela e nota; sem esta linha ela cairia na
% Computer Modern Sans, que não pertence ao conjunto. Uma família só.
\renewcommand{\sfdefault}{\rmdefault}

% DIFERENÇA: mancha de manual, e não a coluna de 468 pt da origem. A4 com
% 22 mm de margem de cada lado, 166 mm úteis. O texto corre em 152 mm e
% os 14 mm da esquerda são o recuo dos títulos: o bastante para o título
% sair do alinhamento do texto, sem deixar coluna vazia ao lado dele. A
% primeira versão reservava 40 mm, e a margem ficava em branco em toda
% página que não fosse título; o documento crescia um terço só por isso.
% `\margemmanual` é quanto os títulos, o cabeçalho e o rodapé avançam
% para a esquerda; vale zero quando o documento traz mancha própria.
% A altura é de 253 mm, a mesma da mancha anterior, porque as tabelas
% deitadas foram dimensionadas para ela.
\newlength{\margemmanual}
\@ifpackageloaded{geometry}{\setlength{\margemmanual}{0pt}}{%
  \setlength{\margemmanual}{14mm}%
  \usepackage[a4paper,left=36mm,right=22mm,top=20mm,bottom=24mm,
              headheight=14pt,headsep=7mm,footskip=12mm]{geometry}}

\usepackage{amsmath}
\usepackage[dvipsnames,table]{xcolor}
\usepackage{graphicx}
\usepackage{tikz}
\usepackage{booktabs}
\usepackage{tabularx}
\usepackage{array}
\usepackage{xltabular}
\usepackage{float}
\usepackage{pdflscape}
\usepackage{rotating}
% Página com o conteúdo girado, mas exibida em retrato no leitor de PDF
\newenvironment{paisagemretrato}{\begingroup\def\PLS@AddRotate##1{}\landscape}{\endlandscape\endgroup}
\usepackage{enumitem}
\usepackage{changepage}
\usepackage[most]{tcolorbox}
\usepackage{caption}
\usepackage{titlesec}
\usepackage{fancyhdr}
\usepackage{lastpage}
\usepackage{csquotes}
\usepackage{siunitx}
% DIFERENÇA: sem `biblatex`. Nenhum documento daqui tem `.bib`; as
% referências são listas escritas à mão.
\usepackage{xurl}
\usepackage[hidelinks,pdfusetitle]{hyperref}

% ---------------------------------------------------------------------
%  Números e unidades (ISO 80000-1)
% ---------------------------------------------------------------------
\sisetup{
  output-decimal-marker = {,},
  group-separator       = {\,},
  group-minimum-digits  = 4,
  exponent-product      = \cdot,
  uncertainty-mode      = separate,
  range-phrase          = {\,a\,},
  list-pair-separator   = {\ e\ },
  list-final-separator  = {\ e\ },
  per-mode              = symbol,
}

% ---------------------------------------------------------------------
%  Parágrafo, listas e flutuantes
% ---------------------------------------------------------------------
% DIFERENÇA: parágrafo em bloco, separado por espaço e sem recuo, que é
% o de manual; a origem recua 1 em.
\setlength{\parindent}{0pt}
\setlength{\parskip}{0.45em plus 0.1em}
\linespread{1.04}
\setlist{itemsep=0.1em,topsep=0.3em,leftmargin=1.4em}
\renewcommand{\arraystretch}{1.15}
\setlength{\tabcolsep}{5pt}
\setlength{\LTpre}{0.8em}
\setlength{\LTpost}{0.4em}

% Legendas no padrão Elsevier: \enquote{Fig. 1.} seguido do texto, alinhado à esquerda;
% \enquote{Tabela 1} numa linha e o título na linha seguinte
\DeclareCaptionLabelSeparator{travessao}{\ — }
% `babel` guarda os nomes por idioma, e o nome da opção varia entre os
% documentos. Entra em cada um que estiver carregado.
\@for\pb@idioma:=brazilian,brazil,portuguese,portuges,american,english\do{%
  \@ifundefined{captions\pb@idioma}{}{%
    \expandafter\addto\csname captions\pb@idioma\endcsname{\renewcommand{\figurename}{Fig.}}}}
\captionsetup{font=footnotesize,labelfont=bf,justification=justified,
              singlelinecheck=false,skip=5pt}
\captionsetup[figure]{position=bottom,labelsep=period}
\captionsetup[table]{position=top,labelsep=newline}
% DIFERENÇA: a origem fixa todo flutuante onde aparece (`H`). Aqui ele
% tenta o lugar e, se não couber no resto da página, sobe para a seguinte
% e o texto continua: com `H`, a tabela que não cabia deixava o fim da
% página em branco.
\floatplacement{table}{!htbp}
\floatplacement{figure}{!htbp}
\renewcommand{\topfraction}{0.9}
\renewcommand{\bottomfraction}{0.8}
\renewcommand{\textfraction}{0.08}
\renewcommand{\floatpagefraction}{0.75}

\newcolumntype{Y}{>{\raggedright\arraybackslash}X}
\newcolumntype{L}[1]{>{\raggedright\arraybackslash}p{#1}}
\newcolumntype{C}[1]{>{\centering\arraybackslash}p{#1}}
\newcolumntype{R}[1]{>{\raggedleft\arraybackslash}p{#1}}

% ---------------------------------------------------------------------
%  Títulos de seção (ISO 2145)
% ---------------------------------------------------------------------
% DIFERENÇA: título de manual. Seção e subseção avançam 14 mm para a
% esquerda do texto e ocupam a largura útil; a seção leva um fio acima.
% Número sem ponto final, três níveis numerados. Os três níveis têm o
% corpo do texto ou pouco mais: o nível se lê pelo fio, pelo recuo e
% pelo itálico do terceiro, e não pelo tamanho.
\setcounter{secnumdepth}{3}
\titleformat{\section}[block]
  {\titlerule[0.8pt]\vspace{0.5ex}\normalfont\large\bfseries\raggedright}{\thesection}{0.8em}{}
\titleformat{\subsection}{\normalfont\normalsize\bfseries\raggedright}{\thesubsection}{0.8em}{}
\titleformat{\subsubsection}{\normalfont\normalsize\bfseries\itshape\raggedright}{\thesubsubsection}{0.8em}{}
\titlespacing*{\section}{-\margemmanual}{1.5em plus 0.4em minus 0.2em}{0.5em plus 0.1em}
\titlespacing*{\subsection}{-\margemmanual}{1.1em plus 0.3em minus 0.2em}{0.3em plus 0.1em}
\titlespacing*{\subsubsection}{0pt}{0.9em plus 0.2em minus 0.2em}{0.2em}

% ---------------------------------------------------------------------
%  Identificação, cabeçalho e rodapé
% ---------------------------------------------------------------------
% Todo campo é do documento, por `\renewcommand` depois deste arquivo.
% Sem isso: o código é o nome do arquivo, a data é a da compilação em
% ISO 8601, e versão e responsável ficam vazios e somem.
\providecommand{\doctitulo}{}
\providecommand{\docsubtitulo}{}
\providecommand{\doccodigo}{\jobname}
\providecommand{\docversao}{}
\providecommand{\docresponsavel}{}
\providecommand{\docdata}{\the\year-\two@digits\month-\two@digits\day}

% Cabeçalho: título à esquerda, seção corrente à direita. Rodapé: código,
% versão e data à esquerda, "página X de Y" à direita. Os dois ocupam a
% largura útil, margem de título incluída.
\newcommand{\pb@identificacao}{%
  \texttt{\doccodigo}%
  \ifx\docversao\@empty\else\enspace\textperiodcentered\enspace versão \docversao\fi
  \enspace\textperiodcentered\enspace\docdata}
\pagestyle{fancy}
\fancyhf{}
\fancyhfoffset[L]{\margemmanual}
\fancyhead[L]{\footnotesize\bfseries\doctitulo}
\fancyhead[R]{\footnotesize\itshape\nouppercase{\rightmark}}
\fancyfoot[L]{\footnotesize\pb@identificacao}
\fancyfoot[R]{\footnotesize\thepage\ de \pageref*{LastPage}}
\renewcommand{\headrulewidth}{0.4pt}
\renewcommand{\footrulewidth}{0.4pt}
\fancypagestyle{plain}{}
% Página de rosto: sem cabeçalho, porque o título já está nela.
\fancypagestyle{rosto}{\fancyhead{}\renewcommand{\headrulewidth}{0pt}}

% ---------------------------------------------------------------------
%  Marca
% ---------------------------------------------------------------------
% DIFERENÇA: a origem não tem marca. Aqui todo documento abre com o
% logotipo Solara, versão para fundo claro. O caminho depende de onde se
% compila: `report/` ou `docs/product/`.
\IfFileExists{../docs/img/solara-logo-on-light.png}%
  {\newcommand{\marcaarquivo}{../docs/img/solara-logo-on-light.png}}%
  {\newcommand{\marcaarquivo}{../img/solara-logo-on-light.png}}
% \marcadoc[largura]: o logotipo tem proporção 7:1, então 30 mm dão
% 4,3 mm de altura.
\newcommand{\marcadoc}[1][30mm]{\includegraphics[width=#1]{\marcaarquivo}}

% ---------------------------------------------------------------------
%  Rosto e histórico de revisões
% ---------------------------------------------------------------------
% \rosto: marca, título, subtítulo e a tabela de identificação. Lê os
% campos `\doc...`; linha de campo vazio não sai.
\newcommand{\pb@linha}[2]{\ifx#2\@empty\else #1 & #2 \\\fi}
\newcommand{\rosto}{%
  \thispagestyle{rosto}%
  \begin{adjustwidth}{-\margemmanual}{0pt}
    \raggedright
    \marcadoc[45mm]\par\vspace{16mm}
    {\fontsize{24}{29}\selectfont\bfseries\doctitulo\par}
    \ifx\docsubtitulo\@empty\else
      \vspace{3mm}{\Large\color{cinza}\docsubtitulo\par}\fi
    \vspace{9mm}{\hrule height 0.8pt}\vspace{4mm}
    {\small\renewcommand{\arraystretch}{1.25}%
     \begin{tabular}{@{}>{\bfseries}p{28mm}p{110mm}@{}}
       Documento & \texttt{\doccodigo} \\
       \pb@linha{Versão}{\docversao}%
       Data & \docdata \\
       \pb@linha{Responsável}{\docresponsavel}%
     \end{tabular}\par}
  \end{adjustwidth}
  \vspace{10mm}}

% Histórico de revisões: uma linha por versão, a mais recente em cima.
%   \begin{historico}
%     0.2 & 2026-10-03 & O que mudou \\
%   \end{historico}
\newenvironment{historico}
  {\par\medskip\noindent\small
   \begin{tabular}{@{}p{16mm}p{24mm}p{\dimexpr\linewidth-40mm-4\tabcolsep\relax}@{}}
   \toprule \textbf{Versão} & \textbf{Data} & \textbf{Alteração} \\ \midrule}
  {\bottomrule\end{tabular}\par\medskip}

% ---------------------------------------------------------------------
%  Avisos (ANSI Z535.6) e nota
% ---------------------------------------------------------------------
% Quatro palavras de sinal, em hierarquia fixa. As três primeiras são de
% dano a pessoa; AVISO é de dano a equipamento ou a dado. Cada mensagem
% diz o risco, como evitar e a consequência. NOTA não é aviso: é
% informação de apoio.
\definecolor{sinalperigo}{HTML}{C8102E}
\definecolor{sinalatencao}{HTML}{E87722}
\definecolor{sinalcuidado}{HTML}{F2C200}
\definecolor{sinalaviso}{HTML}{00539B}
\definecolor{sinalnota}{HTML}{5F5F5F}
% #1 palavra, #2 cor da faixa, #3 cor da letra sobre a faixa
\newtcolorbox{pb@mensagem}[3]{enhanced,breakable,sharp corners,
  boxrule=0pt,frame hidden,borderline west={3pt}{0pt}{#2},
  colback=#2!7,left=9pt,right=7pt,top=5pt,bottom=5pt,
  before skip=0.9em,after skip=0.9em,fontupper=\small,
  before upper={\colorbox{#2}{\textcolor{#3}{\footnotesize\bfseries #1}}\enspace}}
\newenvironment{perigo}{\csname pb@mensagem\endcsname{PERIGO}{sinalperigo}{white}}{\csname endpb@mensagem\endcsname}
\newenvironment{atencao}{\csname pb@mensagem\endcsname{ATENÇÃO}{sinalatencao}{black}}{\csname endpb@mensagem\endcsname}
\newenvironment{cuidado}{\csname pb@mensagem\endcsname{CUIDADO}{sinalcuidado}{black}}{\csname endpb@mensagem\endcsname}
\newenvironment{aviso}{\csname pb@mensagem\endcsname{AVISO}{sinalaviso}{white}}{\csname endpb@mensagem\endcsname}
\newenvironment{nota}{\csname pb@mensagem\endcsname{NOTA}{sinalnota}{white}}{\csname endpb@mensagem\endcsname}

% ---------------------------------------------------------------------
%  Procedimento
% ---------------------------------------------------------------------
% Passos numerados, uma ação por passo, no imperativo. A pré-condição vem
% antes da lista e o resultado esperado depois.
\newlist{passos}{enumerate}{2}
\setlist[passos,1]{label=\textbf{\arabic*.},leftmargin=2.2em,itemsep=0.45em,topsep=0.6em}
\setlist[passos,2]{label=\alph*),leftmargin=1.8em,itemsep=0.2em,topsep=0.2em}
\newcommand{\precondicao}[1]{\par\smallskip\noindent\textbf{Antes de começar.}\enspace#1\par}
\newcommand{\resultado}[1]{\par\smallskip\noindent\textbf{Resultado.}\enspace#1\par}

% Bloco na largura útil inteira, margem de título incluída: para tabela
% ou figura que não cabe nos 152 mm do texto.
\newenvironment{larga}{\begin{adjustwidth}{-\margemmanual}{0pt}}{\end{adjustwidth}}
% O mesmo dentro de `table` ou `figure`, onde `adjustwidth` não alcança:
% uma caixa na largura útil, puxada para a margem. Não quebra página.
% Dentro dela a largura é `\linewidth`, e não `\textwidth`.
\newenvironment{tabelalarga}
  {\noindent\hspace*{-\margemmanual}%
   \begin{minipage}{\dimexpr\linewidth+\margemmanual\relax}\centering}
  {\end{minipage}}

% ---------------------------------------------------------------------
%  Campos e figuras
% ---------------------------------------------------------------------
\definecolor{cinza}{gray}{0.40}

% \figuradoc[largura]{caminho}; sem o arquivo, uma caixa marca a pendência
\newcommand{\figuradoc}[2][\linewidth]{%
  \IfFileExists{#2}%
    {\includegraphics[width=#1,height=0.80\textheight,keepaspectratio]{#2}}%
    {\fbox{\parbox[c][45mm][c]{\dimexpr#1-2\fboxsep-2\fboxrule\relax}%
      {\centering\small\color{cinza}\itshape[figura em preparação]}}}}

\pbrestauraarroba

%
%  FORMA DE MANUAL TÉCNICO. Fonte, números, legendas e pacotes vêm da
%  origem. A forma da página não: é de manual de referência, e não de
%  artigo. Coluna única, títulos recuados para a esquerda do texto,
%  rodapé com identificação e "página X de Y" em toda página, avisos com
%  palavra de sinal (ANSI Z535.6), passos numerados e histórico de
%  revisões (IEC/IEEE 82079-1). O desenho de página segue a classe
%  `refman` do CTAN.
%
%  O documento carrega `babel` ANTES deste arquivo. Documento deitado, que
%  existe para caber tabela larga, carrega também `geometry` antes; aí a
%  margem de manual não entra e os títulos não avançam.
%
%  Motor: pdflatex. Charis SIL, newtxmath e Bera Mono entram como Type 1
%  em T1, e a matemática do newtxmath não passa pelo `unicode-math`.
% =====================================================================
\RequirePackage{iftex}
\RequirePDFTeX
% Este arquivo entra tanto de um `.tex` quanto de dentro do
% `estilo-reporte.sty`, onde o `@` ja e letra. O codigo de categoria e
% devolvido como estava, e nao forcado a "outro" no fim.
\edef\pbrestauraarroba{\catcode`\noexpand\@=\the\catcode`\@\relax}
\makeatletter

% ---------------------------------------------------------------------
%  Fontes, idioma e microtipografia
% ---------------------------------------------------------------------
\usepackage[T1]{fontenc}
\usepackage[utf8]{inputenc}
% DIFERENÇA: o idioma é do documento, porque há documento em inglês.
\@ifpackageloaded{babel}{}{\usepackage[brazilian]{babel}}
\usepackage{CharisSIL}
\usepackage[charter]{newtxmath}
\usepackage[scaled=0.83]{beramono}
\usepackage[protrusion=true,expansion=true,tracking=true]{microtype}
\SetTracking{encoding=*,shape=sc}{40}
% DIFERENÇA: a origem não usa sem serifa. Os documentos daqui pedem
% `\sffamily` em título, tabela e nota; sem esta linha ela cairia na
% Computer Modern Sans, que não pertence ao conjunto. Uma família só.
\renewcommand{\sfdefault}{\rmdefault}

% DIFERENÇA: mancha de manual, e não a coluna de 468 pt da origem. A4 com
% 22 mm de margem de cada lado, 166 mm úteis. O texto corre em 152 mm e
% os 14 mm da esquerda são o recuo dos títulos: o bastante para o título
% sair do alinhamento do texto, sem deixar coluna vazia ao lado dele. A
% primeira versão reservava 40 mm, e a margem ficava em branco em toda
% página que não fosse título; o documento crescia um terço só por isso.
% `\margemmanual` é quanto os títulos, o cabeçalho e o rodapé avançam
% para a esquerda; vale zero quando o documento traz mancha própria.
% A altura é de 253 mm, a mesma da mancha anterior, porque as tabelas
% deitadas foram dimensionadas para ela.
\newlength{\margemmanual}
\@ifpackageloaded{geometry}{\setlength{\margemmanual}{0pt}}{%
  \setlength{\margemmanual}{14mm}%
  \usepackage[a4paper,left=36mm,right=22mm,top=20mm,bottom=24mm,
              headheight=14pt,headsep=7mm,footskip=12mm]{geometry}}

\usepackage{amsmath}
\usepackage[dvipsnames,table]{xcolor}
\usepackage{graphicx}
\usepackage{tikz}
\usepackage{booktabs}
\usepackage{tabularx}
\usepackage{array}
\usepackage{xltabular}
\usepackage{float}
\usepackage{pdflscape}
\usepackage{rotating}
% Página com o conteúdo girado, mas exibida em retrato no leitor de PDF
\newenvironment{paisagemretrato}{\begingroup\def\PLS@AddRotate##1{}\landscape}{\endlandscape\endgroup}
\usepackage{enumitem}
\usepackage{changepage}
\usepackage[most]{tcolorbox}
\usepackage{caption}
\usepackage{titlesec}
\usepackage{fancyhdr}
\usepackage{lastpage}
\usepackage{csquotes}
\usepackage{siunitx}
% DIFERENÇA: sem `biblatex`. Nenhum documento daqui tem `.bib`; as
% referências são listas escritas à mão.
\usepackage{xurl}
\usepackage[hidelinks,pdfusetitle]{hyperref}

% ---------------------------------------------------------------------
%  Números e unidades (ISO 80000-1)
% ---------------------------------------------------------------------
\sisetup{
  output-decimal-marker = {,},
  group-separator       = {\,},
  group-minimum-digits  = 4,
  exponent-product      = \cdot,
  uncertainty-mode      = separate,
  range-phrase          = {\,a\,},
  list-pair-separator   = {\ e\ },
  list-final-separator  = {\ e\ },
  per-mode              = symbol,
}

% ---------------------------------------------------------------------
%  Parágrafo, listas e flutuantes
% ---------------------------------------------------------------------
% DIFERENÇA: parágrafo em bloco, separado por espaço e sem recuo, que é
% o de manual; a origem recua 1 em.
\setlength{\parindent}{0pt}
\setlength{\parskip}{0.45em plus 0.1em}
\linespread{1.04}
\setlist{itemsep=0.1em,topsep=0.3em,leftmargin=1.4em}
\renewcommand{\arraystretch}{1.15}
\setlength{\tabcolsep}{5pt}
\setlength{\LTpre}{0.8em}
\setlength{\LTpost}{0.4em}

% Legendas no padrão Elsevier: \enquote{Fig. 1.} seguido do texto, alinhado à esquerda;
% \enquote{Tabela 1} numa linha e o título na linha seguinte
\DeclareCaptionLabelSeparator{travessao}{\ — }
% `babel` guarda os nomes por idioma, e o nome da opção varia entre os
% documentos. Entra em cada um que estiver carregado.
\@for\pb@idioma:=brazilian,brazil,portuguese,portuges,american,english\do{%
  \@ifundefined{captions\pb@idioma}{}{%
    \expandafter\addto\csname captions\pb@idioma\endcsname{\renewcommand{\figurename}{Fig.}}}}
\captionsetup{font=footnotesize,labelfont=bf,justification=justified,
              singlelinecheck=false,skip=5pt}
\captionsetup[figure]{position=bottom,labelsep=period}
\captionsetup[table]{position=top,labelsep=newline}
% DIFERENÇA: a origem fixa todo flutuante onde aparece (`H`). Aqui ele
% tenta o lugar e, se não couber no resto da página, sobe para a seguinte
% e o texto continua: com `H`, a tabela que não cabia deixava o fim da
% página em branco.
\floatplacement{table}{!htbp}
\floatplacement{figure}{!htbp}
\renewcommand{\topfraction}{0.9}
\renewcommand{\bottomfraction}{0.8}
\renewcommand{\textfraction}{0.08}
\renewcommand{\floatpagefraction}{0.75}

\newcolumntype{Y}{>{\raggedright\arraybackslash}X}
\newcolumntype{L}[1]{>{\raggedright\arraybackslash}p{#1}}
\newcolumntype{C}[1]{>{\centering\arraybackslash}p{#1}}
\newcolumntype{R}[1]{>{\raggedleft\arraybackslash}p{#1}}

% ---------------------------------------------------------------------
%  Títulos de seção (ISO 2145)
% ---------------------------------------------------------------------
% DIFERENÇA: título de manual. Seção e subseção avançam 14 mm para a
% esquerda do texto e ocupam a largura útil; a seção leva um fio acima.
% Número sem ponto final, três níveis numerados. Os três níveis têm o
% corpo do texto ou pouco mais: o nível se lê pelo fio, pelo recuo e
% pelo itálico do terceiro, e não pelo tamanho.
\setcounter{secnumdepth}{3}
\titleformat{\section}[block]
  {\titlerule[0.8pt]\vspace{0.5ex}\normalfont\large\bfseries\raggedright}{\thesection}{0.8em}{}
\titleformat{\subsection}{\normalfont\normalsize\bfseries\raggedright}{\thesubsection}{0.8em}{}
\titleformat{\subsubsection}{\normalfont\normalsize\bfseries\itshape\raggedright}{\thesubsubsection}{0.8em}{}
\titlespacing*{\section}{-\margemmanual}{1.5em plus 0.4em minus 0.2em}{0.5em plus 0.1em}
\titlespacing*{\subsection}{-\margemmanual}{1.1em plus 0.3em minus 0.2em}{0.3em plus 0.1em}
\titlespacing*{\subsubsection}{0pt}{0.9em plus 0.2em minus 0.2em}{0.2em}

% ---------------------------------------------------------------------
%  Identificação, cabeçalho e rodapé
% ---------------------------------------------------------------------
% Todo campo é do documento, por `\renewcommand` depois deste arquivo.
% Sem isso: o código é o nome do arquivo, a data é a da compilação em
% ISO 8601, e versão e responsável ficam vazios e somem.
\providecommand{\doctitulo}{}
\providecommand{\docsubtitulo}{}
\providecommand{\doccodigo}{\jobname}
\providecommand{\docversao}{}
\providecommand{\docresponsavel}{}
\providecommand{\docdata}{\the\year-\two@digits\month-\two@digits\day}

% Cabeçalho: título à esquerda, seção corrente à direita. Rodapé: código,
% versão e data à esquerda, "página X de Y" à direita. Os dois ocupam a
% largura útil, margem de título incluída.
\newcommand{\pb@identificacao}{%
  \texttt{\doccodigo}%
  \ifx\docversao\@empty\else\enspace\textperiodcentered\enspace versão \docversao\fi
  \enspace\textperiodcentered\enspace\docdata}
\pagestyle{fancy}
\fancyhf{}
\fancyhfoffset[L]{\margemmanual}
\fancyhead[L]{\footnotesize\bfseries\doctitulo}
\fancyhead[R]{\footnotesize\itshape\nouppercase{\rightmark}}
\fancyfoot[L]{\footnotesize\pb@identificacao}
\fancyfoot[R]{\footnotesize\thepage\ de \pageref*{LastPage}}
\renewcommand{\headrulewidth}{0.4pt}
\renewcommand{\footrulewidth}{0.4pt}
\fancypagestyle{plain}{}
% Página de rosto: sem cabeçalho, porque o título já está nela.
\fancypagestyle{rosto}{\fancyhead{}\renewcommand{\headrulewidth}{0pt}}

% ---------------------------------------------------------------------
%  Marca
% ---------------------------------------------------------------------
% DIFERENÇA: a origem não tem marca. Aqui todo documento abre com o
% logotipo Solara, versão para fundo claro. O caminho depende de onde se
% compila: `report/` ou `docs/product/`.
\IfFileExists{../docs/img/solara-logo-on-light.png}%
  {\newcommand{\marcaarquivo}{../docs/img/solara-logo-on-light.png}}%
  {\newcommand{\marcaarquivo}{../img/solara-logo-on-light.png}}
% \marcadoc[largura]: o logotipo tem proporção 7:1, então 30 mm dão
% 4,3 mm de altura.
\newcommand{\marcadoc}[1][30mm]{\includegraphics[width=#1]{\marcaarquivo}}

% ---------------------------------------------------------------------
%  Rosto e histórico de revisões
% ---------------------------------------------------------------------
% \rosto: marca, título, subtítulo e a tabela de identificação. Lê os
% campos `\doc...`; linha de campo vazio não sai.
\newcommand{\pb@linha}[2]{\ifx#2\@empty\else #1 & #2 \\\fi}
\newcommand{\rosto}{%
  \thispagestyle{rosto}%
  \begin{adjustwidth}{-\margemmanual}{0pt}
    \raggedright
    \marcadoc[45mm]\par\vspace{16mm}
    {\fontsize{24}{29}\selectfont\bfseries\doctitulo\par}
    \ifx\docsubtitulo\@empty\else
      \vspace{3mm}{\Large\color{cinza}\docsubtitulo\par}\fi
    \vspace{9mm}{\hrule height 0.8pt}\vspace{4mm}
    {\small\renewcommand{\arraystretch}{1.25}%
     \begin{tabular}{@{}>{\bfseries}p{28mm}p{110mm}@{}}
       Documento & \texttt{\doccodigo} \\
       \pb@linha{Versão}{\docversao}%
       Data & \docdata \\
       \pb@linha{Responsável}{\docresponsavel}%
     \end{tabular}\par}
  \end{adjustwidth}
  \vspace{10mm}}

% Histórico de revisões: uma linha por versão, a mais recente em cima.
%   \begin{historico}
%     0.2 & 2026-10-03 & O que mudou \\
%   \end{historico}
\newenvironment{historico}
  {\par\medskip\noindent\small
   \begin{tabular}{@{}p{16mm}p{24mm}p{\dimexpr\linewidth-40mm-4\tabcolsep\relax}@{}}
   \toprule \textbf{Versão} & \textbf{Data} & \textbf{Alteração} \\ \midrule}
  {\bottomrule\end{tabular}\par\medskip}

% ---------------------------------------------------------------------
%  Avisos (ANSI Z535.6) e nota
% ---------------------------------------------------------------------
% Quatro palavras de sinal, em hierarquia fixa. As três primeiras são de
% dano a pessoa; AVISO é de dano a equipamento ou a dado. Cada mensagem
% diz o risco, como evitar e a consequência. NOTA não é aviso: é
% informação de apoio.
\definecolor{sinalperigo}{HTML}{C8102E}
\definecolor{sinalatencao}{HTML}{E87722}
\definecolor{sinalcuidado}{HTML}{F2C200}
\definecolor{sinalaviso}{HTML}{00539B}
\definecolor{sinalnota}{HTML}{5F5F5F}
% #1 palavra, #2 cor da faixa, #3 cor da letra sobre a faixa
\newtcolorbox{pb@mensagem}[3]{enhanced,breakable,sharp corners,
  boxrule=0pt,frame hidden,borderline west={3pt}{0pt}{#2},
  colback=#2!7,left=9pt,right=7pt,top=5pt,bottom=5pt,
  before skip=0.9em,after skip=0.9em,fontupper=\small,
  before upper={\colorbox{#2}{\textcolor{#3}{\footnotesize\bfseries #1}}\enspace}}
\newenvironment{perigo}{\csname pb@mensagem\endcsname{PERIGO}{sinalperigo}{white}}{\csname endpb@mensagem\endcsname}
\newenvironment{atencao}{\csname pb@mensagem\endcsname{ATENÇÃO}{sinalatencao}{black}}{\csname endpb@mensagem\endcsname}
\newenvironment{cuidado}{\csname pb@mensagem\endcsname{CUIDADO}{sinalcuidado}{black}}{\csname endpb@mensagem\endcsname}
\newenvironment{aviso}{\csname pb@mensagem\endcsname{AVISO}{sinalaviso}{white}}{\csname endpb@mensagem\endcsname}
\newenvironment{nota}{\csname pb@mensagem\endcsname{NOTA}{sinalnota}{white}}{\csname endpb@mensagem\endcsname}

% ---------------------------------------------------------------------
%  Procedimento
% ---------------------------------------------------------------------
% Passos numerados, uma ação por passo, no imperativo. A pré-condição vem
% antes da lista e o resultado esperado depois.
\newlist{passos}{enumerate}{2}
\setlist[passos,1]{label=\textbf{\arabic*.},leftmargin=2.2em,itemsep=0.45em,topsep=0.6em}
\setlist[passos,2]{label=\alph*),leftmargin=1.8em,itemsep=0.2em,topsep=0.2em}
\newcommand{\precondicao}[1]{\par\smallskip\noindent\textbf{Antes de começar.}\enspace#1\par}
\newcommand{\resultado}[1]{\par\smallskip\noindent\textbf{Resultado.}\enspace#1\par}

% Bloco na largura útil inteira, margem de título incluída: para tabela
% ou figura que não cabe nos 152 mm do texto.
\newenvironment{larga}{\begin{adjustwidth}{-\margemmanual}{0pt}}{\end{adjustwidth}}
% O mesmo dentro de `table` ou `figure`, onde `adjustwidth` não alcança:
% uma caixa na largura útil, puxada para a margem. Não quebra página.
% Dentro dela a largura é `\linewidth`, e não `\textwidth`.
\newenvironment{tabelalarga}
  {\noindent\hspace*{-\margemmanual}%
   \begin{minipage}{\dimexpr\linewidth+\margemmanual\relax}\centering}
  {\end{minipage}}

% ---------------------------------------------------------------------
%  Campos e figuras
% ---------------------------------------------------------------------
\definecolor{cinza}{gray}{0.40}

% \figuradoc[largura]{caminho}; sem o arquivo, uma caixa marca a pendência
\newcommand{\figuradoc}[2][\linewidth]{%
  \IfFileExists{#2}%
    {\includegraphics[width=#1,height=0.80\textheight,keepaspectratio]{#2}}%
    {\fbox{\parbox[c][45mm][c]{\dimexpr#1-2\fboxsep-2\fboxrule\relax}%
      {\centering\small\color{cinza}\itshape[figura em preparação]}}}}

\pbrestauraarroba

%
%  FORMA DE MANUAL TÉCNICO. Fonte, números, legendas e pacotes vêm da
%  origem. A forma da página não: é de manual de referência, e não de
%  artigo. Coluna única, títulos recuados para a esquerda do texto,
%  rodapé com identificação e "página X de Y" em toda página, avisos com
%  palavra de sinal (ANSI Z535.6), passos numerados e histórico de
%  revisões (IEC/IEEE 82079-1). O desenho de página segue a classe
%  `refman` do CTAN.
%
%  O documento carrega `babel` ANTES deste arquivo. Documento deitado, que
%  existe para caber tabela larga, carrega também `geometry` antes; aí a
%  margem de manual não entra e os títulos não avançam.
%
%  Motor: pdflatex. Charis SIL, newtxmath e Bera Mono entram como Type 1
%  em T1, e a matemática do newtxmath não passa pelo `unicode-math`.
% =====================================================================
\RequirePackage{iftex}
\RequirePDFTeX
% Este arquivo entra tanto de um `.tex` quanto de dentro do
% `estilo-reporte.sty`, onde o `@` ja e letra. O codigo de categoria e
% devolvido como estava, e nao forcado a "outro" no fim.
\edef\pbrestauraarroba{\catcode`\noexpand\@=\the\catcode`\@\relax}
\makeatletter

% ---------------------------------------------------------------------
%  Fontes, idioma e microtipografia
% ---------------------------------------------------------------------
\usepackage[T1]{fontenc}
\usepackage[utf8]{inputenc}
% DIFERENÇA: o idioma é do documento, porque há documento em inglês.
\@ifpackageloaded{babel}{}{\usepackage[brazilian]{babel}}
\usepackage{CharisSIL}
\usepackage[charter]{newtxmath}
\usepackage[scaled=0.83]{beramono}
\usepackage[protrusion=true,expansion=true,tracking=true]{microtype}
\SetTracking{encoding=*,shape=sc}{40}
% DIFERENÇA: a origem não usa sem serifa. Os documentos daqui pedem
% `\sffamily` em título, tabela e nota; sem esta linha ela cairia na
% Computer Modern Sans, que não pertence ao conjunto. Uma família só.
\renewcommand{\sfdefault}{\rmdefault}

% DIFERENÇA: mancha de manual, e não a coluna de 468 pt da origem. A4 com
% 22 mm de margem de cada lado, 166 mm úteis. O texto corre em 152 mm e
% os 14 mm da esquerda são o recuo dos títulos: o bastante para o título
% sair do alinhamento do texto, sem deixar coluna vazia ao lado dele. A
% primeira versão reservava 40 mm, e a margem ficava em branco em toda
% página que não fosse título; o documento crescia um terço só por isso.
% `\margemmanual` é quanto os títulos, o cabeçalho e o rodapé avançam
% para a esquerda; vale zero quando o documento traz mancha própria.
% A altura é de 253 mm, a mesma da mancha anterior, porque as tabelas
% deitadas foram dimensionadas para ela.
\newlength{\margemmanual}
\@ifpackageloaded{geometry}{\setlength{\margemmanual}{0pt}}{%
  \setlength{\margemmanual}{14mm}%
  \usepackage[a4paper,left=36mm,right=22mm,top=20mm,bottom=24mm,
              headheight=14pt,headsep=7mm,footskip=12mm]{geometry}}

\usepackage{amsmath}
\usepackage[dvipsnames,table]{xcolor}
\usepackage{graphicx}
\usepackage{tikz}
\usepackage{booktabs}
\usepackage{tabularx}
\usepackage{array}
\usepackage{xltabular}
\usepackage{float}
\usepackage{pdflscape}
\usepackage{rotating}
% Página com o conteúdo girado, mas exibida em retrato no leitor de PDF
\newenvironment{paisagemretrato}{\begingroup\def\PLS@AddRotate##1{}\landscape}{\endlandscape\endgroup}
\usepackage{enumitem}
\usepackage{changepage}
\usepackage[most]{tcolorbox}
\usepackage{caption}
\usepackage{titlesec}
\usepackage{fancyhdr}
\usepackage{lastpage}
\usepackage{csquotes}
\usepackage{siunitx}
% DIFERENÇA: sem `biblatex`. Nenhum documento daqui tem `.bib`; as
% referências são listas escritas à mão.
\usepackage{xurl}
\usepackage[hidelinks,pdfusetitle]{hyperref}

% ---------------------------------------------------------------------
%  Números e unidades (ISO 80000-1)
% ---------------------------------------------------------------------
\sisetup{
  output-decimal-marker = {,},
  group-separator       = {\,},
  group-minimum-digits  = 4,
  exponent-product      = \cdot,
  uncertainty-mode      = separate,
  range-phrase          = {\,a\,},
  list-pair-separator   = {\ e\ },
  list-final-separator  = {\ e\ },
  per-mode              = symbol,
}

% ---------------------------------------------------------------------
%  Parágrafo, listas e flutuantes
% ---------------------------------------------------------------------
% DIFERENÇA: parágrafo em bloco, separado por espaço e sem recuo, que é
% o de manual; a origem recua 1 em.
\setlength{\parindent}{0pt}
\setlength{\parskip}{0.45em plus 0.1em}
\linespread{1.04}
\setlist{itemsep=0.1em,topsep=0.3em,leftmargin=1.4em}
\renewcommand{\arraystretch}{1.15}
\setlength{\tabcolsep}{5pt}
\setlength{\LTpre}{0.8em}
\setlength{\LTpost}{0.4em}

% Legendas no padrão Elsevier: \enquote{Fig. 1.} seguido do texto, alinhado à esquerda;
% \enquote{Tabela 1} numa linha e o título na linha seguinte
\DeclareCaptionLabelSeparator{travessao}{\ — }
% `babel` guarda os nomes por idioma, e o nome da opção varia entre os
% documentos. Entra em cada um que estiver carregado.
\@for\pb@idioma:=brazilian,brazil,portuguese,portuges,american,english\do{%
  \@ifundefined{captions\pb@idioma}{}{%
    \expandafter\addto\csname captions\pb@idioma\endcsname{\renewcommand{\figurename}{Fig.}}}}
\captionsetup{font=footnotesize,labelfont=bf,justification=justified,
              singlelinecheck=false,skip=5pt}
\captionsetup[figure]{position=bottom,labelsep=period}
\captionsetup[table]{position=top,labelsep=newline}
% DIFERENÇA: a origem fixa todo flutuante onde aparece (`H`). Aqui ele
% tenta o lugar e, se não couber no resto da página, sobe para a seguinte
% e o texto continua: com `H`, a tabela que não cabia deixava o fim da
% página em branco.
\floatplacement{table}{!htbp}
\floatplacement{figure}{!htbp}
\renewcommand{\topfraction}{0.9}
\renewcommand{\bottomfraction}{0.8}
\renewcommand{\textfraction}{0.08}
\renewcommand{\floatpagefraction}{0.75}

\newcolumntype{Y}{>{\raggedright\arraybackslash}X}
\newcolumntype{L}[1]{>{\raggedright\arraybackslash}p{#1}}
\newcolumntype{C}[1]{>{\centering\arraybackslash}p{#1}}
\newcolumntype{R}[1]{>{\raggedleft\arraybackslash}p{#1}}

% ---------------------------------------------------------------------
%  Títulos de seção (ISO 2145)
% ---------------------------------------------------------------------
% DIFERENÇA: título de manual. Seção e subseção avançam 14 mm para a
% esquerda do texto e ocupam a largura útil; a seção leva um fio acima.
% Número sem ponto final, três níveis numerados. Os três níveis têm o
% corpo do texto ou pouco mais: o nível se lê pelo fio, pelo recuo e
% pelo itálico do terceiro, e não pelo tamanho.
\setcounter{secnumdepth}{3}
\titleformat{\section}[block]
  {\titlerule[0.8pt]\vspace{0.5ex}\normalfont\large\bfseries\raggedright}{\thesection}{0.8em}{}
\titleformat{\subsection}{\normalfont\normalsize\bfseries\raggedright}{\thesubsection}{0.8em}{}
\titleformat{\subsubsection}{\normalfont\normalsize\bfseries\itshape\raggedright}{\thesubsubsection}{0.8em}{}
\titlespacing*{\section}{-\margemmanual}{1.5em plus 0.4em minus 0.2em}{0.5em plus 0.1em}
\titlespacing*{\subsection}{-\margemmanual}{1.1em plus 0.3em minus 0.2em}{0.3em plus 0.1em}
\titlespacing*{\subsubsection}{0pt}{0.9em plus 0.2em minus 0.2em}{0.2em}

% ---------------------------------------------------------------------
%  Identificação, cabeçalho e rodapé
% ---------------------------------------------------------------------
% Todo campo é do documento, por `\renewcommand` depois deste arquivo.
% Sem isso: o código é o nome do arquivo, a data é a da compilação em
% ISO 8601, e versão e responsável ficam vazios e somem.
\providecommand{\doctitulo}{}
\providecommand{\docsubtitulo}{}
\providecommand{\doccodigo}{\jobname}
\providecommand{\docversao}{}
\providecommand{\docresponsavel}{}
\providecommand{\docdata}{\the\year-\two@digits\month-\two@digits\day}

% Cabeçalho: título à esquerda, seção corrente à direita. Rodapé: código,
% versão e data à esquerda, "página X de Y" à direita. Os dois ocupam a
% largura útil, margem de título incluída.
\newcommand{\pb@identificacao}{%
  \texttt{\doccodigo}%
  \ifx\docversao\@empty\else\enspace\textperiodcentered\enspace versão \docversao\fi
  \enspace\textperiodcentered\enspace\docdata}
\pagestyle{fancy}
\fancyhf{}
\fancyhfoffset[L]{\margemmanual}
\fancyhead[L]{\footnotesize\bfseries\doctitulo}
\fancyhead[R]{\footnotesize\itshape\nouppercase{\rightmark}}
\fancyfoot[L]{\footnotesize\pb@identificacao}
\fancyfoot[R]{\footnotesize\thepage\ de \pageref*{LastPage}}
\renewcommand{\headrulewidth}{0.4pt}
\renewcommand{\footrulewidth}{0.4pt}
\fancypagestyle{plain}{}
% Página de rosto: sem cabeçalho, porque o título já está nela.
\fancypagestyle{rosto}{\fancyhead{}\renewcommand{\headrulewidth}{0pt}}

% ---------------------------------------------------------------------
%  Marca
% ---------------------------------------------------------------------
% DIFERENÇA: a origem não tem marca. Aqui todo documento abre com o
% logotipo Solara, versão para fundo claro. O caminho depende de onde se
% compila: `report/` ou `docs/product/`.
\IfFileExists{../docs/img/solara-logo-on-light.png}%
  {\newcommand{\marcaarquivo}{../docs/img/solara-logo-on-light.png}}%
  {\newcommand{\marcaarquivo}{../img/solara-logo-on-light.png}}
% \marcadoc[largura]: o logotipo tem proporção 7:1, então 30 mm dão
% 4,3 mm de altura.
\newcommand{\marcadoc}[1][30mm]{\includegraphics[width=#1]{\marcaarquivo}}

% ---------------------------------------------------------------------
%  Rosto e histórico de revisões
% ---------------------------------------------------------------------
% \rosto: marca, título, subtítulo e a tabela de identificação. Lê os
% campos `\doc...`; linha de campo vazio não sai.
\newcommand{\pb@linha}[2]{\ifx#2\@empty\else #1 & #2 \\\fi}
\newcommand{\rosto}{%
  \thispagestyle{rosto}%
  \begin{adjustwidth}{-\margemmanual}{0pt}
    \raggedright
    \marcadoc[45mm]\par\vspace{16mm}
    {\fontsize{24}{29}\selectfont\bfseries\doctitulo\par}
    \ifx\docsubtitulo\@empty\else
      \vspace{3mm}{\Large\color{cinza}\docsubtitulo\par}\fi
    \vspace{9mm}{\hrule height 0.8pt}\vspace{4mm}
    {\small\renewcommand{\arraystretch}{1.25}%
     \begin{tabular}{@{}>{\bfseries}p{28mm}p{110mm}@{}}
       Documento & \texttt{\doccodigo} \\
       \pb@linha{Versão}{\docversao}%
       Data & \docdata \\
       \pb@linha{Responsável}{\docresponsavel}%
     \end{tabular}\par}
  \end{adjustwidth}
  \vspace{10mm}}

% Histórico de revisões: uma linha por versão, a mais recente em cima.
%   \begin{historico}
%     0.2 & 2026-10-03 & O que mudou \\
%   \end{historico}
\newenvironment{historico}
  {\par\medskip\noindent\small
   \begin{tabular}{@{}p{16mm}p{24mm}p{\dimexpr\linewidth-40mm-4\tabcolsep\relax}@{}}
   \toprule \textbf{Versão} & \textbf{Data} & \textbf{Alteração} \\ \midrule}
  {\bottomrule\end{tabular}\par\medskip}

% ---------------------------------------------------------------------
%  Avisos (ANSI Z535.6) e nota
% ---------------------------------------------------------------------
% Quatro palavras de sinal, em hierarquia fixa. As três primeiras são de
% dano a pessoa; AVISO é de dano a equipamento ou a dado. Cada mensagem
% diz o risco, como evitar e a consequência. NOTA não é aviso: é
% informação de apoio.
\definecolor{sinalperigo}{HTML}{C8102E}
\definecolor{sinalatencao}{HTML}{E87722}
\definecolor{sinalcuidado}{HTML}{F2C200}
\definecolor{sinalaviso}{HTML}{00539B}
\definecolor{sinalnota}{HTML}{5F5F5F}
% #1 palavra, #2 cor da faixa, #3 cor da letra sobre a faixa
\newtcolorbox{pb@mensagem}[3]{enhanced,breakable,sharp corners,
  boxrule=0pt,frame hidden,borderline west={3pt}{0pt}{#2},
  colback=#2!7,left=9pt,right=7pt,top=5pt,bottom=5pt,
  before skip=0.9em,after skip=0.9em,fontupper=\small,
  before upper={\colorbox{#2}{\textcolor{#3}{\footnotesize\bfseries #1}}\enspace}}
\newenvironment{perigo}{\csname pb@mensagem\endcsname{PERIGO}{sinalperigo}{white}}{\csname endpb@mensagem\endcsname}
\newenvironment{atencao}{\csname pb@mensagem\endcsname{ATENÇÃO}{sinalatencao}{black}}{\csname endpb@mensagem\endcsname}
\newenvironment{cuidado}{\csname pb@mensagem\endcsname{CUIDADO}{sinalcuidado}{black}}{\csname endpb@mensagem\endcsname}
\newenvironment{aviso}{\csname pb@mensagem\endcsname{AVISO}{sinalaviso}{white}}{\csname endpb@mensagem\endcsname}
\newenvironment{nota}{\csname pb@mensagem\endcsname{NOTA}{sinalnota}{white}}{\csname endpb@mensagem\endcsname}

% ---------------------------------------------------------------------
%  Procedimento
% ---------------------------------------------------------------------
% Passos numerados, uma ação por passo, no imperativo. A pré-condição vem
% antes da lista e o resultado esperado depois.
\newlist{passos}{enumerate}{2}
\setlist[passos,1]{label=\textbf{\arabic*.},leftmargin=2.2em,itemsep=0.45em,topsep=0.6em}
\setlist[passos,2]{label=\alph*),leftmargin=1.8em,itemsep=0.2em,topsep=0.2em}
\newcommand{\precondicao}[1]{\par\smallskip\noindent\textbf{Antes de começar.}\enspace#1\par}
\newcommand{\resultado}[1]{\par\smallskip\noindent\textbf{Resultado.}\enspace#1\par}

% Bloco na largura útil inteira, margem de título incluída: para tabela
% ou figura que não cabe nos 152 mm do texto.
\newenvironment{larga}{\begin{adjustwidth}{-\margemmanual}{0pt}}{\end{adjustwidth}}
% O mesmo dentro de `table` ou `figure`, onde `adjustwidth` não alcança:
% uma caixa na largura útil, puxada para a margem. Não quebra página.
% Dentro dela a largura é `\linewidth`, e não `\textwidth`.
\newenvironment{tabelalarga}
  {\noindent\hspace*{-\margemmanual}%
   \begin{minipage}{\dimexpr\linewidth+\margemmanual\relax}\centering}
  {\end{minipage}}

% ---------------------------------------------------------------------
%  Campos e figuras
% ---------------------------------------------------------------------
\definecolor{cinza}{gray}{0.40}

% \figuradoc[largura]{caminho}; sem o arquivo, uma caixa marca a pendência
\newcommand{\figuradoc}[2][\linewidth]{%
  \IfFileExists{#2}%
    {\includegraphics[width=#1,height=0.80\textheight,keepaspectratio]{#2}}%
    {\fbox{\parbox[c][45mm][c]{\dimexpr#1-2\fboxsep-2\fboxrule\relax}%
      {\centering\small\color{cinza}\itshape[figura em preparação]}}}}

\pbrestauraarroba

%
%  FORMA DE MANUAL TÉCNICO. Fonte, números, legendas e pacotes vêm da
%  origem. A forma da página não: é de manual de referência, e não de
%  artigo. Coluna única, títulos recuados para a esquerda do texto,
%  rodapé com identificação e "página X de Y" em toda página, avisos com
%  palavra de sinal (ANSI Z535.6), passos numerados e histórico de
%  revisões (IEC/IEEE 82079-1). O desenho de página segue a classe
%  `refman` do CTAN.
%
%  O documento carrega `babel` ANTES deste arquivo. Documento deitado, que
%  existe para caber tabela larga, carrega também `geometry` antes; aí a
%  margem de manual não entra e os títulos não avançam.
%
%  Motor: pdflatex. Charis SIL, newtxmath e Bera Mono entram como Type 1
%  em T1, e a matemática do newtxmath não passa pelo `unicode-math`.
% =====================================================================
\RequirePackage{iftex}
\RequirePDFTeX
% Este arquivo entra tanto de um `.tex` quanto de dentro do
% `estilo-reporte.sty`, onde o `@` ja e letra. O codigo de categoria e
% devolvido como estava, e nao forcado a "outro" no fim.
\edef\pbrestauraarroba{\catcode`\noexpand\@=\the\catcode`\@\relax}
\makeatletter

% ---------------------------------------------------------------------
%  Fontes, idioma e microtipografia
% ---------------------------------------------------------------------
\usepackage[T1]{fontenc}
\usepackage[utf8]{inputenc}
% DIFERENÇA: o idioma é do documento, porque há documento em inglês.
\@ifpackageloaded{babel}{}{\usepackage[brazilian]{babel}}
\usepackage{CharisSIL}
\usepackage[charter]{newtxmath}
\usepackage[scaled=0.83]{beramono}
\usepackage[protrusion=true,expansion=true,tracking=true]{microtype}
\SetTracking{encoding=*,shape=sc}{40}
% DIFERENÇA: a origem não usa sem serifa. Os documentos daqui pedem
% `\sffamily` em título, tabela e nota; sem esta linha ela cairia na
% Computer Modern Sans, que não pertence ao conjunto. Uma família só.
\renewcommand{\sfdefault}{\rmdefault}

% DIFERENÇA: mancha de manual, e não a coluna de 468 pt da origem. A4 com
% 22 mm de margem de cada lado, 166 mm úteis. O texto corre em 152 mm e
% os 14 mm da esquerda são o recuo dos títulos: o bastante para o título
% sair do alinhamento do texto, sem deixar coluna vazia ao lado dele. A
% primeira versão reservava 40 mm, e a margem ficava em branco em toda
% página que não fosse título; o documento crescia um terço só por isso.
% `\margemmanual` é quanto os títulos, o cabeçalho e o rodapé avançam
% para a esquerda; vale zero quando o documento traz mancha própria.
% A altura é de 253 mm, a mesma da mancha anterior, porque as tabelas
% deitadas foram dimensionadas para ela.
\newlength{\margemmanual}
\@ifpackageloaded{geometry}{\setlength{\margemmanual}{0pt}}{%
  \setlength{\margemmanual}{14mm}%
  \usepackage[a4paper,left=36mm,right=22mm,top=20mm,bottom=24mm,
              headheight=14pt,headsep=7mm,footskip=12mm]{geometry}}

\usepackage{amsmath}
\usepackage[dvipsnames,table]{xcolor}
\usepackage{graphicx}
\usepackage{tikz}
\usepackage{booktabs}
\usepackage{tabularx}
\usepackage{array}
\usepackage{xltabular}
\usepackage{float}
\usepackage{pdflscape}
\usepackage{rotating}
% Página com o conteúdo girado, mas exibida em retrato no leitor de PDF
\newenvironment{paisagemretrato}{\begingroup\def\PLS@AddRotate##1{}\landscape}{\endlandscape\endgroup}
\usepackage{enumitem}
\usepackage{changepage}
\usepackage[most]{tcolorbox}
\usepackage{caption}
\usepackage{titlesec}
\usepackage{fancyhdr}
\usepackage{lastpage}
\usepackage{csquotes}
\usepackage{siunitx}
% DIFERENÇA: sem `biblatex`. Nenhum documento daqui tem `.bib`; as
% referências são listas escritas à mão.
\usepackage{xurl}
\usepackage[hidelinks,pdfusetitle]{hyperref}

% ---------------------------------------------------------------------
%  Números e unidades (ISO 80000-1)
% ---------------------------------------------------------------------
\sisetup{
  output-decimal-marker = {,},
  group-separator       = {\,},
  group-minimum-digits  = 4,
  exponent-product      = \cdot,
  uncertainty-mode      = separate,
  range-phrase          = {\,a\,},
  list-pair-separator   = {\ e\ },
  list-final-separator  = {\ e\ },
  per-mode              = symbol,
}

% ---------------------------------------------------------------------
%  Parágrafo, listas e flutuantes
% ---------------------------------------------------------------------
% DIFERENÇA: parágrafo em bloco, separado por espaço e sem recuo, que é
% o de manual; a origem recua 1 em.
\setlength{\parindent}{0pt}
\setlength{\parskip}{0.45em plus 0.1em}
\linespread{1.04}
\setlist{itemsep=0.1em,topsep=0.3em,leftmargin=1.4em}
\renewcommand{\arraystretch}{1.15}
\setlength{\tabcolsep}{5pt}
\setlength{\LTpre}{0.8em}
\setlength{\LTpost}{0.4em}

% Legendas no padrão Elsevier: \enquote{Fig. 1.} seguido do texto, alinhado à esquerda;
% \enquote{Tabela 1} numa linha e o título na linha seguinte
\DeclareCaptionLabelSeparator{travessao}{\ — }
% `babel` guarda os nomes por idioma, e o nome da opção varia entre os
% documentos. Entra em cada um que estiver carregado.
\@for\pb@idioma:=brazilian,brazil,portuguese,portuges,american,english\do{%
  \@ifundefined{captions\pb@idioma}{}{%
    \expandafter\addto\csname captions\pb@idioma\endcsname{\renewcommand{\figurename}{Fig.}}}}
\captionsetup{font=footnotesize,labelfont=bf,justification=justified,
              singlelinecheck=false,skip=5pt}
\captionsetup[figure]{position=bottom,labelsep=period}
\captionsetup[table]{position=top,labelsep=newline}
% DIFERENÇA: a origem fixa todo flutuante onde aparece (`H`). Aqui ele
% tenta o lugar e, se não couber no resto da página, sobe para a seguinte
% e o texto continua: com `H`, a tabela que não cabia deixava o fim da
% página em branco.
\floatplacement{table}{!htbp}
\floatplacement{figure}{!htbp}
\renewcommand{\topfraction}{0.9}
\renewcommand{\bottomfraction}{0.8}
\renewcommand{\textfraction}{0.08}
\renewcommand{\floatpagefraction}{0.75}

\newcolumntype{Y}{>{\raggedright\arraybackslash}X}
\newcolumntype{L}[1]{>{\raggedright\arraybackslash}p{#1}}
\newcolumntype{C}[1]{>{\centering\arraybackslash}p{#1}}
\newcolumntype{R}[1]{>{\raggedleft\arraybackslash}p{#1}}

% ---------------------------------------------------------------------
%  Títulos de seção (ISO 2145)
% ---------------------------------------------------------------------
% DIFERENÇA: título de manual. Seção e subseção avançam 14 mm para a
% esquerda do texto e ocupam a largura útil; a seção leva um fio acima.
% Número sem ponto final, três níveis numerados. Os três níveis têm o
% corpo do texto ou pouco mais: o nível se lê pelo fio, pelo recuo e
% pelo itálico do terceiro, e não pelo tamanho.
\setcounter{secnumdepth}{3}
\titleformat{\section}[block]
  {\titlerule[0.8pt]\vspace{0.5ex}\normalfont\large\bfseries\raggedright}{\thesection}{0.8em}{}
\titleformat{\subsection}{\normalfont\normalsize\bfseries\raggedright}{\thesubsection}{0.8em}{}
\titleformat{\subsubsection}{\normalfont\normalsize\bfseries\itshape\raggedright}{\thesubsubsection}{0.8em}{}
\titlespacing*{\section}{-\margemmanual}{1.5em plus 0.4em minus 0.2em}{0.5em plus 0.1em}
\titlespacing*{\subsection}{-\margemmanual}{1.1em plus 0.3em minus 0.2em}{0.3em plus 0.1em}
\titlespacing*{\subsubsection}{0pt}{0.9em plus 0.2em minus 0.2em}{0.2em}

% ---------------------------------------------------------------------
%  Identificação, cabeçalho e rodapé
% ---------------------------------------------------------------------
% Todo campo é do documento, por `\renewcommand` depois deste arquivo.
% Sem isso: o código é o nome do arquivo, a data é a da compilação em
% ISO 8601, e versão e responsável ficam vazios e somem.
\providecommand{\doctitulo}{}
\providecommand{\docsubtitulo}{}
\providecommand{\doccodigo}{\jobname}
\providecommand{\docversao}{}
\providecommand{\docresponsavel}{}
\providecommand{\docdata}{\the\year-\two@digits\month-\two@digits\day}

% Cabeçalho: título à esquerda, seção corrente à direita. Rodapé: código,
% versão e data à esquerda, "página X de Y" à direita. Os dois ocupam a
% largura útil, margem de título incluída.
\newcommand{\pb@identificacao}{%
  \texttt{\doccodigo}%
  \ifx\docversao\@empty\else\enspace\textperiodcentered\enspace versão \docversao\fi
  \enspace\textperiodcentered\enspace\docdata}
\pagestyle{fancy}
\fancyhf{}
\fancyhfoffset[L]{\margemmanual}
\fancyhead[L]{\footnotesize\bfseries\doctitulo}
\fancyhead[R]{\footnotesize\itshape\nouppercase{\rightmark}}
\fancyfoot[L]{\footnotesize\pb@identificacao}
\fancyfoot[R]{\footnotesize\thepage\ de \pageref*{LastPage}}
\renewcommand{\headrulewidth}{0.4pt}
\renewcommand{\footrulewidth}{0.4pt}
\fancypagestyle{plain}{}
% Página de rosto: sem cabeçalho, porque o título já está nela.
\fancypagestyle{rosto}{\fancyhead{}\renewcommand{\headrulewidth}{0pt}}

% ---------------------------------------------------------------------
%  Marca
% ---------------------------------------------------------------------
% DIFERENÇA: a origem não tem marca. Aqui todo documento abre com o
% logotipo Solara, versão para fundo claro. O caminho depende de onde se
% compila: `report/` ou `docs/product/`.
\IfFileExists{../docs/img/solara-logo-on-light.png}%
  {\newcommand{\marcaarquivo}{../docs/img/solara-logo-on-light.png}}%
  {\newcommand{\marcaarquivo}{../img/solara-logo-on-light.png}}
% \marcadoc[largura]: o logotipo tem proporção 7:1, então 30 mm dão
% 4,3 mm de altura.
\newcommand{\marcadoc}[1][30mm]{\includegraphics[width=#1]{\marcaarquivo}}

% ---------------------------------------------------------------------
%  Rosto e histórico de revisões
% ---------------------------------------------------------------------
% \rosto: marca, título, subtítulo e a tabela de identificação. Lê os
% campos `\doc...`; linha de campo vazio não sai.
\newcommand{\pb@linha}[2]{\ifx#2\@empty\else #1 & #2 \\\fi}
\newcommand{\rosto}{%
  \thispagestyle{rosto}%
  \begin{adjustwidth}{-\margemmanual}{0pt}
    \raggedright
    \marcadoc[45mm]\par\vspace{16mm}
    {\fontsize{24}{29}\selectfont\bfseries\doctitulo\par}
    \ifx\docsubtitulo\@empty\else
      \vspace{3mm}{\Large\color{cinza}\docsubtitulo\par}\fi
    \vspace{9mm}{\hrule height 0.8pt}\vspace{4mm}
    {\small\renewcommand{\arraystretch}{1.25}%
     \begin{tabular}{@{}>{\bfseries}p{28mm}p{110mm}@{}}
       Documento & \texttt{\doccodigo} \\
       \pb@linha{Versão}{\docversao}%
       Data & \docdata \\
       \pb@linha{Responsável}{\docresponsavel}%
     \end{tabular}\par}
  \end{adjustwidth}
  \vspace{10mm}}

% Histórico de revisões: uma linha por versão, a mais recente em cima.
%   \begin{historico}
%     0.2 & 2026-10-03 & O que mudou \\
%   \end{historico}
\newenvironment{historico}
  {\par\medskip\noindent\small
   \begin{tabular}{@{}p{16mm}p{24mm}p{\dimexpr\linewidth-40mm-4\tabcolsep\relax}@{}}
   \toprule \textbf{Versão} & \textbf{Data} & \textbf{Alteração} \\ \midrule}
  {\bottomrule\end{tabular}\par\medskip}

% ---------------------------------------------------------------------
%  Avisos (ANSI Z535.6) e nota
% ---------------------------------------------------------------------
% Quatro palavras de sinal, em hierarquia fixa. As três primeiras são de
% dano a pessoa; AVISO é de dano a equipamento ou a dado. Cada mensagem
% diz o risco, como evitar e a consequência. NOTA não é aviso: é
% informação de apoio.
\definecolor{sinalperigo}{HTML}{C8102E}
\definecolor{sinalatencao}{HTML}{E87722}
\definecolor{sinalcuidado}{HTML}{F2C200}
\definecolor{sinalaviso}{HTML}{00539B}
\definecolor{sinalnota}{HTML}{5F5F5F}
% #1 palavra, #2 cor da faixa, #3 cor da letra sobre a faixa
\newtcolorbox{pb@mensagem}[3]{enhanced,breakable,sharp corners,
  boxrule=0pt,frame hidden,borderline west={3pt}{0pt}{#2},
  colback=#2!7,left=9pt,right=7pt,top=5pt,bottom=5pt,
  before skip=0.9em,after skip=0.9em,fontupper=\small,
  before upper={\colorbox{#2}{\textcolor{#3}{\footnotesize\bfseries #1}}\enspace}}
\newenvironment{perigo}{\csname pb@mensagem\endcsname{PERIGO}{sinalperigo}{white}}{\csname endpb@mensagem\endcsname}
\newenvironment{atencao}{\csname pb@mensagem\endcsname{ATENÇÃO}{sinalatencao}{black}}{\csname endpb@mensagem\endcsname}
\newenvironment{cuidado}{\csname pb@mensagem\endcsname{CUIDADO}{sinalcuidado}{black}}{\csname endpb@mensagem\endcsname}
\newenvironment{aviso}{\csname pb@mensagem\endcsname{AVISO}{sinalaviso}{white}}{\csname endpb@mensagem\endcsname}
\newenvironment{nota}{\csname pb@mensagem\endcsname{NOTA}{sinalnota}{white}}{\csname endpb@mensagem\endcsname}

% ---------------------------------------------------------------------
%  Procedimento
% ---------------------------------------------------------------------
% Passos numerados, uma ação por passo, no imperativo. A pré-condição vem
% antes da lista e o resultado esperado depois.
\newlist{passos}{enumerate}{2}
\setlist[passos,1]{label=\textbf{\arabic*.},leftmargin=2.2em,itemsep=0.45em,topsep=0.6em}
\setlist[passos,2]{label=\alph*),leftmargin=1.8em,itemsep=0.2em,topsep=0.2em}
\newcommand{\precondicao}[1]{\par\smallskip\noindent\textbf{Antes de começar.}\enspace#1\par}
\newcommand{\resultado}[1]{\par\smallskip\noindent\textbf{Resultado.}\enspace#1\par}

% Bloco na largura útil inteira, margem de título incluída: para tabela
% ou figura que não cabe nos 152 mm do texto.
\newenvironment{larga}{\begin{adjustwidth}{-\margemmanual}{0pt}}{\end{adjustwidth}}
% O mesmo dentro de `table` ou `figure`, onde `adjustwidth` não alcança:
% uma caixa na largura útil, puxada para a margem. Não quebra página.
% Dentro dela a largura é `\linewidth`, e não `\textwidth`.
\newenvironment{tabelalarga}
  {\noindent\hspace*{-\margemmanual}%
   \begin{minipage}{\dimexpr\linewidth+\margemmanual\relax}\centering}
  {\end{minipage}}

% ---------------------------------------------------------------------
%  Campos e figuras
% ---------------------------------------------------------------------
\definecolor{cinza}{gray}{0.40}

% \figuradoc[largura]{caminho}; sem o arquivo, uma caixa marca a pendência
\newcommand{\figuradoc}[2][\linewidth]{%
  \IfFileExists{#2}%
    {\includegraphics[width=#1,height=0.80\textheight,keepaspectratio]{#2}}%
    {\fbox{\parbox[c][45mm][c]{\dimexpr#1-2\fboxsep-2\fboxrule\relax}%
      {\centering\small\color{cinza}\itshape[figura em preparação]}}}}

\pbrestauraarroba
